\documentclass[11pt]{article}
\usepackage[a4paper,margin=1in]{geometry}
\usepackage{times}
\usepackage{amsmath,amssymb}
\usepackage{hyperref}
\usepackage{enumitem}
\hypersetup{hidelinks}

\begin{document}

\begin{center}
\Large\textbf{Statement on Prior Publication and Differences}
\end{center}

\vspace{1em}

\noindent\textbf{Present submission:} Beyond English-Only GRPO: A
Multi-Seed Empirical Study at Sub-3B Scale

\noindent\textbf{Author:} Vu Dang (Independent Researcher, Vietnam)

\noindent\textbf{Date:} 11 May 2026

\vspace{1em}\hrule\vspace{1em}

\subsection*{Prior publication disclosed}

A preliminary single-seed version of this work (referred to as
\textit{v1}) was self-archived as a preprint on Zenodo on 7 May 2026.
The preprint has not undergone peer review.

\medskip

\noindent\textbf{Citation of v1:}

\begin{quote}
Vu Dang. ``Beyond English-Only GRPO: Training Language and Auxiliary
Reward as Implicit Regularizers in Sub-3B Math Reasoning,'' Zenodo,
7~May 2026. DOI: \href{https://doi.org/10.5281/zenodo.20061328}%
{10.5281/zenodo.20061328}.
\end{quote}

The accompanying open-source software repository
(\href{https://github.com/vudang4494/xling-grpo-sub3b}%
{github.com/vudang4494/xling-grpo-sub3b}, Apache-2.0) is shared
between the v1 preprint and the present submission; the repository
contains code, configurations, model checkpoints, and evaluation JSONs
for both versions.

\vspace{1em}

\subsection*{How the present submission (v2) differs from v1}

The submitted manuscript substantially extends the v1 preprint with new
experimental content, methodological additions, and revised conclusions.
The principal differences are:

\begin{enumerate}[itemsep=4pt]

\item \textbf{Multi-seed verification (new experiments).}
v1 reported single-seed results (\texttt{seed=42}) for three training
arms. The present submission expands to three random seeds
(\texttt{seed} $\in \{42, 123, 7\}$) per arm. This reveals that v1's
key claims (e.g., $+7.5$\,pp lift on AMC-23, $+13.3$\,pp recovery on
AIME-2024 via $R_5$) fall within seed variance ($\sigma = 11.3$\,pp
on AMC-23) and require revision. The v2 manuscript reports multi-seed
mean$\,\pm\,$standard deviation and bootstrap 95\,\% confidence
intervals for every quantitative claim.

\item \textbf{Mechanism ablation (new arm).}
v2 adds Arm A4 (constant-bias reward $R_{\text{const}}=1.0$, single
seed) to test the reward-magnitude mechanism hypothesis proposed in v1.
The ablation partially refutes a pure reward-magnitude explanation of
the $R_5$ effect on AIME-2024 maj@8: A4 reverts to base level
($33.3$\,\%) while A3 with the actual $R_5$ reward achieves
$37.8 \pm 1.9$\,\%. This negative result is new in v2 and is reported
honestly as the mechanism remaining an open question.

\item \textbf{Multilingual evaluation (new benchmark).}
v2 adds a 10-language MGSM evaluation across all training arms plus the
untrained base, comprising 80 new evaluation JSONs (8 conditions
$\times$ 10 languages). The result is a negative finding: all training
conditions converge to within $\pm 0.5$\,pp of the untrained base mean,
indicating that training-language choice does \emph{not} produce
cross-lingual transfer at this scale. v1 did not include any MGSM
evaluation.

\item \textbf{Reproducibility infrastructure (new methodology).}
v2 introduces (a)~a SHA256-checksummed manifest covering all 124+
training and evaluation artifacts, (b)~bootstrap 95\,\% confidence
intervals for all numerical claims, (c)~a 13-section reproducibility
audit checklist, and (d)~a per-component verification document
(\texttt{VERIFICATION.md}). v1 reported numbers without these
verification guarantees.

\item \textbf{Reframed contributions (revised narrative).}
v1 framed the work as a positive ``cross-lingual breakthrough'' with
auxiliary reward as a regularization mechanism. v2 reframes around
three orthogonal findings: (i) positive --- A3 best on AIME-2024 maj@8;
(ii) methodological --- multi-seed instability of vanilla EN GRPO;
(iii) negative --- no cross-lingual transfer on MGSM. The mechanism
hypothesis is presented as an open question rather than a confirmed
claim.

\item \textbf{Compute-cost transparency.}
v1 contained references to specific monetary compute amounts (e.g.,
``\$20 single-GPU LoRA budget''). v2 removes all such references for
neutrality, describing hardware in physical terms only
(``one NVIDIA A100 80\,GB GPU'').

\end{enumerate}

\vspace{1em}

\subsection*{Status of v1}

The v1 Zenodo preprint remains publicly accessible at its original DOI
and is cited within the v2 manuscript as
\texttt{[dang\_v1\_preprint\_2026]}. All v1 single-seed numbers are
explicitly marked as superseded by the multi-seed results in v2. The
author intends to upload the v2 manuscript as a new version of the
existing Zenodo DOI upon acceptance, with the v1 version remaining
visible for historical reference.

\vspace{1em}

\subsection*{No other prior publication}

This manuscript has not been submitted to any peer-reviewed venue,
conference, or journal other than IEEE Access. No portions of it have
been presented at any conference or workshop. There are no other
related publications or preprints by the author covering this work.

\vspace{2em}\hrule\vspace{1em}

\noindent Signed:

\medskip

\noindent Vu Dang \\
\href{mailto:vu.dh4494@gmail.com}{vu.dh4494@gmail.com} \\
11 May 2026

\end{document}
